\documentclass{training}
% Filigrane PAR MODULE : \uxfond (PNG A4 dilue 30 %). Cadre noir vectoriel commun.
% Proportions original 744 px : marge 20/744, trait 1/744. Exterieur BLANC.
% Cadre via eso-pic +\put (coordonnees absolues) — pas de remember picture (decalage).
\usepackage{eso-pic}

\providecommand{\uxfond}{../../resources/fond/norme_fond_paliers_30.png}

\newlength{\uxmargin}
\newlength{\uxthick}
\newlength{\uxinnerw}
\newlength{\uxinnerh}
\setlength{\uxmargin}{0.02688\paperwidth}% 20/744
\setlength{\uxthick}{0.001344\paperwidth}% 1/744
\setlength{\uxinnerw}{\dimexpr\paperwidth-2\uxmargin-2\uxthick\relax}
\setlength{\uxinnerh}{\dimexpr\paperheight-2\uxmargin-2\uxthick\relax}

% Contenu strictement a l'interieur du trait (+ jour confortable).
\geometry{%
  includeheadfoot,
  left=22mm,
  right=22mm,
  top=20mm,
  bottom=18mm,
  headheight=30pt,
  headsep=5mm,
  footskip=12mm%
}

\AddToShipoutPictureBG{%
  \AtPageLowerLeft{%
    % Fond blanc page entiere
    \put(0,0){\color{white}\rule{\paperwidth}{\paperheight}}%
    % Filigrane module a l'interieur du cadre
    \put(\LenToUnit{\dimexpr\uxmargin+\uxthick\relax},\LenToUnit{\dimexpr\uxmargin+\uxthick\relax}){%
      \includegraphics[width=\uxinnerw,height=\uxinnerh]{\uxfond}}%
    % Cadre noir : bas, haut, gauche, droite
    \put(\LenToUnit{\uxmargin},\LenToUnit{\uxmargin}){%
      \color{black}\rule{\dimexpr\paperwidth-2\uxmargin\relax}{\uxthick}}%
    \put(\LenToUnit{\uxmargin},\LenToUnit{\dimexpr\paperheight-\uxmargin-\uxthick\relax}){%
      \color{black}\rule{\dimexpr\paperwidth-2\uxmargin\relax}{\uxthick}}%
    \put(\LenToUnit{\uxmargin},\LenToUnit{\uxmargin}){%
      \color{black}\rule{\uxthick}{\dimexpr\paperheight-2\uxmargin\relax}}%
    \put(\LenToUnit{\dimexpr\paperwidth-\uxmargin-\uxthick\relax},\LenToUnit{\uxmargin}){%
      \color{black}\rule{\uxthick}{\dimexpr\paperheight-2\uxmargin\relax}}%
  }%
}

\makeatletter
\fancyhf{}
\fancyheadoffset{0pt}
\renewcommand{\headrulewidth}{0pt}
\lhead{\parbox[t]{0.56\textwidth}{\RaggedRight\footnotesize\sloppy\@title}}
\rhead{\parbox[t]{0.36\textwidth}{\RaggedLeft\footnotesize\sloppy\@subtitle}}
\cfoot{\thepage}
% Coupures : underscore et slash dans les noms de fichiers / consignes
\DeclareRobustCommand{\_}{\textunderscore\allowbreak}
\makeatother

% Evite les depassements de longues lignes (Github, Name_FirstName_...)
\sloppy
\emergencystretch=3em
\hbadness=10000

% Sommaire type 42 : chapitres seulement
\setcounter{tocdepth}{0}
\renewcommand{\cftchappresnum}{}
\renewcommand{\cftchapaftersnum}{}
\setlength{\cftchapnumwidth}{2.0em}
\setlength{\cftbeforechapskip}{0.35em}
\cftsetindents{chapter}{0em}{2.0em}


\begin{document}

\schoollogo{../../resources/jeunecrack.png}
\title{Colors and moodboard}
\subtitle{Piscine UX/UI --- UI}
\member{Piscine UX/UI}{https://github.com/01-edu/UX-UI}
\summary{Color, type, a moodboard, and five interface challenges.}
\version{1.0}
\prerequisite{The previous module of this piscine, or none if this is your first.}
\maketitle

\newpage
\tableofcontents
\newpage


\chapter{Instructions}

The foreword is not an exercise. It will not be graded. The exercises start at the next chapter, from the easiest contract to the hardest. A later exercise is not taken into account if an earlier one does not hold.

Only what is in the turn-in directory of an exercise is looked at. Do not add files that the exercise does not name. The name of a deliverable follows the pattern in the exercise: \texttt{Name\_FirstName\_DeliverableName\_Date\_VersionNumber}.

You will not be given the gesture. The examples in an exercise sometimes carry a constraint that the sentence does not repeat. Read them. Your other references are the manual of the tool you picked, the person sitting next to you, and the links printed in the exercise.

The order is the order of the quest. Skipping a step because a later screen looks more interesting produces a screen nobody can explain.

\subsection*{About the pictures}

The mark on the cover is the project logo. The picture in the foreword is the illustration of this module. The pale ground behind the text is the page watermark. None of them is a hint.

\subsection*{About machines that write}

You may use a tool to look up a term, the way you would use a dictionary. You may not replace an interview, a test, or a file with a generated text and sign it. If you cannot point at the person, the screen, or the line you claim, the exercise is not done.


\chapter{Foreword}

Josef Albers spent a book showing that a color is not a chip, it is what the neighbor does to it (\emph{Interaction of Color}, Yale University Press, 1963). Johannes Itten had already made the same point with paper at the Bauhaus. The challenges that follow are timed drawings. Albers does not paint them for you.

\begin{center}
\includegraphics[width=0.78\textwidth,height=0.32\textheight,keepaspectratio]{../../resources/intros/UI-00-colors-and-moodboard.jpg}\\[0.4em]
{\footnotesize Picture: Art Projects MKCL KF / CC BY-SA 4.0}
\end{center}

Recommended watching. The videos are context. They are not the deliverable.
\begin{itemize}
\item UI - Introduction \& Quest n°1, 01Edu \url{https://www.youtube.com/watch?v=UbE6WMeuvQg&list=PLHyAJ_GrRtf-_aBft2jQL4-QAunn2bLvu}
\end{itemize}

\chapter{Exercise 00: UI I - Ex 1 - Find the colors}
\exnumber{00}
\extitle{UI I - Ex 1 - Find the colors}
\exfiles{ex00/Figma file and PDF, named Name\_FirstName\_DeliverableName\_Date\_VersionNumber}
\exauthorize{Those named in the exercise}
\makeheaderfiles

\textbf{Context:}

This first UI quest is a series of exercises to practice colors, moodboards, typography and replicate screens on Figma!

\textbf{Instructions:}

Go to these 7 websites and make the color palettes out of them (3 to 4 main colors).

\textbf{Deliverable:}

\begin{itemize}
\item A Figma file including screenshots of the main page, URL and extracted colors next to the screenshots.
\item Upload the Figma file on Github, labeled as follows: "Name\_FirstName\_DeliverableName\_Date\_VersionNumber". For example, it can be labeled as "Doe\_John\_MidFiPrototype\_05242024\_V1".
\item Add a title within the file.
\item Pay attention to the global aspect of the file. It must be clear, simple and easy to read. You can get inspiration on canva.com but don't overload your design with too much details!
\end{itemize}

\textbf{Tips:}

\begin{itemize}
\item You can use Figma to extract the colors you need and export your file to PDF.
\item Use the Color Pipette to extract the colors!
\item You can show the main colors with filled squares or circles next to the screenshots. 
\item Mind the aesthetics of your delivrable! 1 page per site should be good.
\end{itemize}

\begin{itemize}
\item https://www.komoot.com/ \url{https://www.komoot.com/} (App)
\item https://www.airtable.com/ \url{https://www.airtable.com/}
\item https://www.decathlon.com/collections/kids-gift-ideas \url{https://www.decathlon.com/collections/kids-gift-ideas}
\item https://www.jrpass.com/ \url{https://www.jrpass.com/}
\item https://www.thesiff.com/ \url{https://www.thesiff.com/}
\item https://brave.com/es/ \url{https://brave.com/es/}
\item https://www.nhp.gov.in/ \url{https://www.nhp.gov.in/}
\end{itemize}

\textbf{Resources:}

\begin{itemize}
\item Video tutorial on Figma \url{https://www.youtube.com/watch?v=FTFaQWZBqQ8}
\item Beginning graphic design \url{https://edu.gcfglobal.org/en/beginning-graphic-design/}
\item Gareth graphic studio \url{https://www.youtube.com/c/GarethDavidStudio}
\end{itemize}


\chapter{Exercise 01: UI I - Ex 2 - Color codes}
\exnumber{01}
\extitle{UI I - Ex 2 - Color codes}
\exfiles{ex01/Figma file, named Name\_FirstName\_DeliverableName\_Date\_VersionNumber}
\exauthorize{Those named in the exercise}
\makeheaderfiles

\textbf{Instructions:}

Go to these 7 websites and make the color palettes out of them (3 to 4 main colors).

\textbf{Deliverable:}

\begin{itemize}
\item A Figma file including screenshots of the main page, URL, extracted colors and the matching color codes.
\item Upload the Figma file on Github, labeled as follows: "Name\_FirstName\_DeliverableName\_Date\_VersionNumber". For example, it can be labeled as "Doe\_John\_MidFiPrototype\_05242024\_V1".
\item Add a title within the file.
\item Pay attention to the global aspect of the file. It must be clear, simple and easy to read. You can get inspiration on canva.com but don't overload your design with too much details!
\end{itemize}

\textbf{Tips:}

\begin{itemize}
\item You can use the document from the previous exercise to add the color codes on Figma.
\item Use the Color Pipette to detect the color codes!
\end{itemize}

\textbf{Resources:}

\begin{itemize}
\item Beginning graphic design \url{https://edu.gcfglobal.org/en/beginning-graphic-design/}
\item Gareth graphic studio \url{https://www.youtube.com/c/GarethDavidStudio}
\end{itemize}

\begin{itemize}
\item https://www.komoot.com/ \url{https://www.komoot.com/} (App)
\item https://www.airtable.com/ \url{https://www.airtable.com/}
\item https://www.decathlon.com/collections/kids-gift-ideas \url{https://www.decathlon.com/collections/kids-gift-ideas}
\item https://www.jrpass.com/ \url{https://www.jrpass.com/}
\item https://www.thesiff.com/ \url{https://www.thesiff.com/}
\item https://brave.com/es/ \url{https://brave.com/es/}
\item https://www.nhp.gov.in/ \url{https://www.nhp.gov.in/}
\end{itemize}


\chapter{Exercise 02: UI I - Ex 3 - Color combinations}
\exnumber{02}
\extitle{UI I - Ex 3 - Color combinations}
\exfiles{ex02/Figma file, named Name\_FirstName\_DeliverableName\_Date\_VersionNumber}
\exauthorize{Those named in the exercise}
\makeheaderfiles

\textbf{Instructions:}

Let's browse!

Find 3 digital products (websites or apps) that use each kind of color palette combination:

\begin{itemize}
\item Monochromatic
\item Analogous
\item Complementary
\item Split complementary
\item Triadic
\item Double complementary
\end{itemize}

\textbf{Deliverable:}

\begin{itemize}
\item A Figma file including for each digital product (websites or apps)
\begin{itemize}
\item screenshots of the main page
\item the matching URL
\item the main colors and
\item the kind of color combination.
\end{itemize}
\item Upload the Figma file on Github, labeled as follows: "Name\_FirstName\_DeliverableName\_Date\_VersionNumber". For example, it can be labeled as "Doe\_John\_MidFiPrototype\_05242024\_V1".
\item Add a title within the file.
\item Pay attention to the global aspect of the file. It must be clear, simple and easy to read. You can get inspiration on canva.com but don't overload your design with too much details!
\end{itemize}

\textbf{Tips:}

\begin{itemize}
\item You can find HSB (Hue, Saturation, Brightness) as well as HSV (Hue, Saturation, Value) depending of the source.
\end{itemize}

\textbf{Resources:}

\begin{itemize}
\item https://coloursandmaterials.wordpress.com/2014/09/16/colour-system/ \url{https://coloursandmaterials.wordpress.com/2014/09/16/colour-system/}
\item Color wheel combinations
\end{itemize}

\textbf{Monochromatic examples:}

\begin{itemize}
\item https://www.benefitcosmetics.com/ \url{https://www.benefitcosmetics.com/} 
\item http://www.jwilouvres.com.au/ \url{http://www.jwilouvres.com.au/} 
\item http://blank.com.pt/ \url{http://blank.com.pt/}
\end{itemize}

\textbf{Analogous examples:}

\begin{itemize}
\item https://www.natuurenmilieu.nl/ \url{https://www.natuurenmilieu.nl/} 
\item https://campesinorum.com/ \url{https://campesinorum.com/} 
\item https://www.lobster.es/ \url{https://www.lobster.es/}
\end{itemize}

\textbf{Complementary examples:}

\begin{itemize}
\item https://www.rolandgarros.com/ \url{https://www.rolandgarros.com/fr-fr} 
\item https://squilla.io/ \url{https://squilla.io/} 
\item https://gomim.com/en/ \url{https://gomim.com/en/}
\end{itemize}

\textbf{Split complementary examples:}

\begin{itemize}
\item https://flythenest.io/en/ \url{https://flythenest.io/en/} 
\item https://www.kineuphorics.com/ \url{https://www.kineuphorics.com/} 
\item https://www.ma-tea.de/ \url{https://www.ma-tea.de/}
\end{itemize}

\textbf{Triadic examples}

\begin{itemize}
\item https://normalnow.com/ \url{https://normalnow.com/} 
\item https://brainbakery.com/en/ \url{https://brainbakery.com/en/} 
\item https://www.barkbox.com/ \url{https://www.barkbox.com/}
\end{itemize}

\textbf{Dual complementary or Tetradic examples}

\begin{itemize}
\item https://www.airtable.com/ \url{https://www.airtable.com/} 
\item https://calendar.google.com/ \url{https://calendar.google.com/} 
\item https://www.ebay.com/ \url{https://www.ebay.com/}
\end{itemize}


\chapter{Exercise 03: UI I - Ex 4 - Moodboard}
\exnumber{03}
\extitle{UI I - Ex 4 - Moodboard}
\exfiles{ex03/Figma file, named Name\_FirstName\_DeliverableName\_Date\_VersionNumber}
\exauthorize{Those named in the exercise}
\makeheaderfiles

\textbf{Instructions:}

Create 4 different moodboards based on these 4 briefs and extract a color palette.

\textbf{Briefs:}

\begin{itemize}
\item 1: A physics teacher that wants to make his website catchy for 12 year-olds.
\item 2: A Sri Lankan wedding planner website that wants to breathe sobriety and modernity.
\item 3: An e-commerce platform that sells West African wax and wants to feel more luxurious in its spirit to stand out.
\item 4: An NGO that helps parents of kids with autism, wants to spread joy and hope for a fundraiser.
\end{itemize}

\textbf{Deliverables:}

\begin{itemize}
\item A Figma file per moodboard
\item On each moodboard, several pictures (5 to 12) that are graphically consistent in terms of colors
\item The main colors extracted and shown on the side
\item One color code per color (RGB, Hex or HSB)
\item Upload the documents on Github, in a zip folder named "ProjectTitle\_Name\_FirstName", with all project deliverables as follows: "Name\_FirstName\_DeliverableName\_Date\_VersionNumber".  For example, the first deliverable can be named as "Doe\_John\_MidFiPrototype\_05242024\_V1".
\item Add a title within all written documents.
\item Pay attention to the global aspect of the deliverables. They must be clear, simple and easy to read. You can get inspiration on canva.com but don't overload your design with too much details!
\end{itemize}

\textbf{Tips:}

\begin{itemize}
\item Use freely-usable images! \url{https://buffer.com/library/free-images/}
\item Examples of moodboards (in addition, the moodboard should include the color codes): 
\item https://lovestylecomunicacion.com/wp-content/uploads/2017/09/A.png \url{https://lovestylecomunicacion.com/wp-content/uploads/2017/09/A.png} 
\item https://caseperlatesta.com/wp-content/uploads/2015/09/Moodboard23\_grey\_peach-1024x1024.jpg \url{https://caseperlatesta.com/wp-content/uploads/2015/09/Moodboard23_grey_peach-1024x1024.jpg} 
\item https://i.pinimg.com/564x/ee/9d/d6/ee9dd6fbb82b0e5b4dd2924580b765dd.jpg \url{https://i.pinimg.com/564x/ee/9d/d6/ee9dd6fbb82b0e5b4dd2924580b765dd.jpg}
\end{itemize}


\chapter{Exercise 04: UI I - Ex 5 - Typography}
\exnumber{04}
\extitle{UI I - Ex 5 - Typography}
\exfiles{ex04/Figma file, named Name\_FirstName\_DeliverableName\_Date\_VersionNumber}
\exauthorize{Those named in the exercise}
\makeheaderfiles

\textbf{Instructions:}

Go to these websites and detect the typographies used.

\begin{itemize}
\item https://unric.org/no \url{https://unric.org/no}
\item https://www.strongbowcider.com/bug-exterminator-brisbane \url{https://www.strongbowcider.com/bug-exterminator-brisbane}
\item http://www.bbqgrillsandfire.com/ \url{http://www.bbqgrillsandfire.com/}
\item https://pt.wikipedia.org/wiki/Mario\_Botta \url{https://pt.wikipedia.org/wiki/Mario_Botta}
\item https://texas-private-investigator.com/ \url{https://texas-private-investigator.com/}
\item https://bumble.com/ \url{https://bumble.com/}
\item https://adopte1poule.fr/ \url{https://adopte1poule.fr/}
\item http://slaveryfootprint.org/ \url{http://slaveryfootprint.org/}
\item https://www.tomorrowland.com/fr/festival/bienvenue \url{https://www.tomorrowland.com/fr/festival/bienvenue}
\item https://artisanscanada.com/ \url{https://artisanscanada.com/}
\end{itemize}

\textbf{Deliverables:}

\begin{itemize}
\item A Figma file showing all the typographies used in each website or app.
\item Upload the Figma file on Github, labeled as follows: "Name\_FirstName\_DeliverableName\_Date\_VersionNumber". For example, it can be labeled as "Doe\_John\_MidFiPrototype\_05242024\_V1".
\item Add a title within the file.
\item Pay attention to the global aspect of the file. It must be clear, simple and easy to read. You can get inspiration on canva.com but don't overload your design with too much details!
\end{itemize}

\textbf{Tips:}

\begin{itemize}
\item Mind the presentation!
\item Feel free to browse Google Chrome options to detect typography ;-)
\item Google Fonts \url{https://fonts.google.com/} is another great friend to have by your side
\end{itemize}

\textbf{Resources:}

\begin{itemize}
\item What is Typography \url{https://www.renderforest.com/blog/what-is-typography}
\item Why is Typography important? \url{https://careerfoundry.com/en/blog/ui-design/beginners-guide-to-typography/}
\item Beginning graphic design: Typography \url{https://edu.gcfglobal.org/en/beginning-graphic-design/typography/1/}
\item Beginning graphic design \url{https://edu.gcfglobal.org/en/beginning-graphic-design/}
\item Gareth graphic studio \url{https://www.youtube.com/c/GarethDavidStudio}
\end{itemize}


\chapter{Exercise 05: UI I - Ex 6 - Typography Moodboard}
\exnumber{05}
\extitle{UI I - Ex 6 - Typography Moodboard}
\exfiles{ex05/Figma file, named Name\_FirstName\_DeliverableName\_Date\_VersionNumber}
\exauthorize{Those named in the exercise}
\makeheaderfiles

\textbf{Instructions:}

Find consistent typography for your moodboards from Ex. 4.

\textbf{Deliverables:}

\begin{itemize}
\item A Figma file with 4 moodboards from Ex. 4 and typography:
\begin{itemize}
\item 1 for a title
\item 1 for a button
\item 1 for a description
\item 1 for the text
\end{itemize}
\item Upload the Figma file on Github, labeled as follows: "Name\_FirstName\_DeliverableName\_Date\_VersionNumber". For example, it can be labeled as "Doe\_John\_MidFiPrototype\_05242024\_V1".
\item Add a title within the file.
\item Pay attention to the global aspect of the file. It must be clear, simple and easy to read. You can get inspiration on canva.com but don't overload your design with too much details!
\end{itemize}

\textbf{Resources:}

\begin{itemize}
\item Font \url{https://en.wikipedia.org/wiki/Font}
\item Glyph \url{https://en.wikipedia.org/wiki/Glyph}
\item Instagram account: You are typography \url{https://www.instagram.com/youaretypography/}
\item Instagram account: Typespire \url{https://www.instagram.com/typespire/}
\item Instagram account: ux\_ui.world \url{https://www.instagram.com/ui_ux.world/}
\item Instagram account: ui.netwrk \url{https://www.instagram.com/ui.netwrk/}
\item Instagram account: uxuidailytips \url{https://www.instagram.com/uiuxdailytips/}
\end{itemize}

\textbf{Tips:}

\begin{itemize}
\item Bold, italic, normal, light, etc. are fonts, not typography.
\item Google Fonts \url{https://fonts.google.com/} is another great friend to have by your side
\item Example of a caption here \url{https://www.figma.com/file/SbQMPjfLOzg92g8YZFlrRh/UI-I---Ex-5}
\end{itemize}

\textbf{Resources:}

\begin{itemize}
\item What is Typography \url{https://www.renderforest.com/blog/what-is-typography}
\item Why is Typography important? \url{https://careerfoundry.com/en/blog/ui-design/beginners-guide-to-typography/}
\item Beginning graphic design: Typography \url{https://edu.gcfglobal.org/en/beginning-graphic-design/typography/1/}
\item Beginning graphic design \url{https://edu.gcfglobal.org/en/beginning-graphic-design/}
\item Gareth graphic studio \url{https://www.youtube.com/c/GarethDavidStudio}
\end{itemize}


\chapter{Exercise 06: UI I - Ex 7 - UI Challenge - Timer}
\exnumber{06}
\extitle{UI I - Ex 7 - UI Challenge - Timer}
\exfiles{ex06/Figma file, named Name\_FirstName\_DeliverableName\_Date\_VersionNumber}
\exauthorize{Those named in the exercise}
\makeheaderfiles

\textbf{Instructions:}

Copy this image on Figma.
Timer

\textbf{Deliverable:}

\begin{itemize}
\item A Figma file
\item Please make sure you use:
\begin{itemize}
\item 2 half circles to draw the circle around the timer
\item Gradient to fill the main circle
\item Groups
\item Named layers
\end{itemize}
\item Upload the Figma file on Github, labeled as follows: "Name\_FirstName\_DeliverableName\_Date\_VersionNumber". For example, it can be labeled as "Doe\_John\_MidFiPrototype\_05242024\_V1".
\item Add a title within the file.
\item Pay attention to the global aspect of the file. It must be clear, simple and easy to read. You can get inspiration on canva.com but don't overload your design with too much details!
\end{itemize}

\textbf{Tips:}

\begin{itemize}
\item The font is RAJDANHI
\item The frame is Android 360 x 640px
\end{itemize}


\chapter{Exercise 07: UI I - Ex 8 - UI Challenge - Watches}
\exnumber{07}
\extitle{UI I - Ex 8 - UI Challenge - Watches}
\exfiles{ex07/Figma file, named Name\_FirstName\_DeliverableName\_Date\_VersionNumber}
\exauthorize{Those named in the exercise}
\makeheaderfiles

\textbf{Instructions:}

Copy these images on Figma

Screen1
Screen2
Screen3
Screen4

\textbf{Deliverable:}

\begin{itemize}
\item A Figma file with a copy of the image with Figma elements
\item Make sure:
\begin{itemize}
\item You use symbols and variants
\item All your layers are named and organized
\item You use overlays
\item You use masks to change the shape of the imported image
\end{itemize}
\item Upload the Figma file on Github, labeled as follows: "Name\_FirstName\_DeliverableName\_Date\_VersionNumber". For example, it can be labeled as "Doe\_John\_MidFiPrototype\_05242024\_V1".
\item Add a title within the file.
\item Pay attention to the global aspect of the file. It must be clear, simple and easy to read. You can get inspiration on canva.com but don't overload your design with too much details!
\end{itemize}

\textbf{Tips:}

\begin{itemize}
\item Fonts: SF Compact for the iOS elements,  Futura STD Condensed for the Nike signs
\item The frame for each iWatch screen is watch 44 mm 184 x 224
\item Figma article about Symbols and Variants \url{https://help.figma.com/hc/en-us/articles/360056440594-Create-and-use-variants}
\item You can use existing Libraries by browsing the Figma community resources
\end{itemize}


\chapter{Exercise 08: UI I - Ex 9 - UI Challenge - Flight}
\exnumber{08}
\extitle{UI I - Ex 9 - UI Challenge - Flight}
\exfiles{ex08/Figma file, named Name\_FirstName\_DeliverableName\_Date\_VersionNumber}
\exauthorize{Those named in the exercise}
\makeheaderfiles

\textbf{Instructions}:

Copy this image on Figma

Flight

\textbf{Deliverable:}

\begin{itemize}
\item A Figma file with a copy of the image with Figma elements
\item Make sure:
\begin{itemize}
\item You use symbols and variants
\item All you layers are names and organized
\item You use overlays
\item You use masks to change the shape of an imported image (Amsterdam, Maldives, New Delhi, Sydney icons)
\end{itemize}
\item Upload the Figma file on Github, labeled as follows: "Name\_FirstName\_DeliverableName\_Date\_VersionNumber". For example, it can be labeled as "Doe\_John\_MidFiPrototype\_05242024\_V1".
\item Add a title within the file.
\item Pay attention to the global aspect of the file. It must be clear, simple and easy to read. You can get inspiration on canva.com but don't overload your design with too much details!
\end{itemize}

\textbf{Tips:}

\begin{itemize}
\item Font: Roboto Bold, Normal, Black with 0 to 7\% space between letters
\item The frame is desktop 1440 x 1048 px
\item Figma article about Symbols and Variants \url{https://help.figma.com/hc/en-us/articles/360056440594-Create-and-use-variants}
\end{itemize}


\chapter{Exercise 09: UI I - Ex 10 - UI Challenge - Spotify}
\exnumber{09}
\extitle{UI I - Ex 10 - UI Challenge - Spotify}
\exfiles{ex09/Figma file, named Name\_FirstName\_DeliverableName\_Date\_VersionNumber}
\exauthorize{Those named in the exercise}
\makeheaderfiles

\textbf{Instructions:}

Copy this image on Figma
Spotify

\textbf{Deliverable:}

\begin{itemize}
\item A Figma file with a copy of the image with Figma elements
\item Please make sure you use:
\begin{itemize}
\item Linear color background
\item Groups
\item Named and organized layers
\item Apple design system (head bar)
\item Symbols and variants for the songs (1 page for the Symbols)
\end{itemize}
\item Upload the Figma file on Github, labeled as follows: "Name\_FirstName\_DeliverableName\_Date\_VersionNumber". For example, it can be labeled as "Doe\_John\_MidFiPrototype\_05242024\_V1".
\item Add a title within the file.
\item Pay attention to the global aspect of the file. It must be clear, simple and easy to read. You can get inspiration on canva.com but don't overload your design with too much details!
\end{itemize}

\textbf{Tips:}

\begin{itemize}
\item The font is GOTHAM
\item The frame is iPhone 13 Pro Max 428 x 926 px
\item Figma article about Symbols and Variants \url{https://help.figma.com/hc/en-us/articles/360056440594-Create-and-use-variants}
\end{itemize}


\chapter{Exercise 10: UI I - Ex 11 - UI Challenge - Bitcoin Dashboard}
\exnumber{10}
\extitle{UI I - Ex 11 - UI Challenge - Bitcoin Dashboard}
\exfiles{ex10/Figma file, named Name\_FirstName\_DeliverableName\_Date\_VersionNumber}
\exauthorize{Those named in the exercise}
\makeheaderfiles

\textbf{Instructions:}

Copy this image on Figma
Dashboard

\textbf{Deliverable:}

\begin{itemize}
\item A Figma file with a copy of the image with Figma elements
\item Make sure:
\begin{itemize}
\item You use symbols and variants
\item All your layers are named and organized
\item You use overlays
\item You use masks to change the shape of an imported image
\item You use linear color function
\end{itemize}
\item Upload the Figma file on Github, labeled as follows: "Name\_FirstName\_DeliverableName\_Date\_VersionNumber". For example, it can be labeled as "Doe\_John\_MidFiPrototype\_05242024\_V1".
\item Add a title within the file.
\item Pay attention to the global aspect of the file. It must be clear, simple and easy to read. You can get inspiration on canva.com but don't overload your design with too much details!
\end{itemize}

\textbf{Tips:}

\begin{itemize}
\item Fonts: Avenir
\item The frame is a desktop 1440 x 1024
\item Figma article about Symbols and Variants \url{https://help.figma.com/hc/en-us/articles/360056440594-Create-and-use-variants}
\item You can use existing Libraries by browsing the Figma community resources
\end{itemize}


\end{document}
